Hidden hunger is a major health concern and is also known as micronutrient deficiency. It is a form of undernutrition resulting from the consumption of energy-dense diets that lack essential vitamins and minerals. Micronutrient deficiencies are not restricted to economically poor countries; they also affect high-income regions such as Europe due to changes in dietary patterns driven by the consumption of processed foods \autocite{1ofori}. \par

The consequences of malnutrition and hidden hunger are severe. Annually, an estimated 1.1 million out of 3.1 million infant deaths are attributed to micronutrient deficiencies \autocites{2c}{largescale}. Less intense forms of malnutrition can cause minor issues like dizziness but can also lead to loss of vision if the malnutrition is advanced. The adverse effect of malnutrition and hidden hunger has forced several countries to address this worldwide problem. International movements, such as the “Scaling Up Nutrition Movement”, but also government programs, seek to fight malnutrition \autocites{4hidden}{5unnes}.  \par

There are four primary interventions that can be differentiated and are addressed to tackle malnutrition and hidden hunger: dietary diversification, food supplementation, food fortification and biofortification \autocites{6biofort}{7plants}. \par

Food fortification involves adding specific nutrients to foods to enhance their nutritional value and help consumers meet the recommended daily allowance (RDA) for those nutrients \autocite{6biofort}. Particularly in developed countries, fortifying foods with nutrients such as iron, iodine and vitamins has greatly reduced the prevalence of diseases associated with nutrient deficiencies \autocite{8forti}. In Germany, for example, salt is iodized to prevent iodine-related disorders, including goiter \autocite{iodine}, and pregnant women are advised to increase their intake of folic acid in order to avoid anemia. These examples illustrate the importance of fortification strategies and their potential public health impact. Hence, it is important to understand every aspect of micronutrient fortification, since the aforementioned problems can be highly individual and require deeper knowledge of the mechanisms that occur during food preparation as well as pre- and post-fortification. \par

Pre-fortification, also referred to as biofortification, increases the nutritional value of crops during their growth. This effect is achieved through various methods such as agronomic practices, conventional plant breeding or genetic engineering \autocite{1ofori}. Biofortification of Vitamin A in sweet potatoes or Golden Rice leads to higher values of Vitamin A in the finished products \autocite{1ofori}. Additionally maize hybrids are being biofortified, to achieve the breeding objective of 15µg/g Provitamin-A Carotenoids, since traditional maize hybrids contain only low amounts (< 2 µg/g) \autocite{dutta}. Being performed only once, prefortification is a highly sustainable method of increasing nutritional value and thus can be viewed as a great long term solution. In contrast, during food preparation there are several mechanisms that hinder the complete preservation of micronutrients. For example, it has been shown that cooking fortified rice in excess water increases micronutrient loss, similar to unfortified rice \autocite{10rice}. 

In the final analysis, post-fortification governs the shelf-life and nutrient stability of the food vehicle and its fortification substance. The consequent implementation of this practice has been used in projects such as ”the fortified biscuit” by the world food program and therefore mirrors the necessity of this practice \autocite{wfp}. \par

To our knowledge, studies on stabilizing micronutrients post-fortification receive less attention in literature than other aspects of fortification. Thus, food producers still rely on overage methods if they want to ensure sufficient amounts of micronutrients to combat hidden hunger. This approach is cost-increasing, hinders new ventures and should therefore be minimized, especially for economically weaker countries. This study aims to summarize recent findings on the topic, to provide greater clarity for food producers who need to ensure that their products contain minimum levels of nutrients in a cost-efficient manner. 
